%% 
%% Template article for cas-sc documentclass for 
%% double column output.

\documentclass[a4paper,fleqn]{cas-sc}

% If the frontmatter runs over more than one page
% Use the longmktitle option.

%\documentclass[a4paper,fleqn,longmktitle]{cas-sc}

%\usepackage[numbers]{natbib}
%\usepackage[authoryear]{natbib}
\usepackage[numbers]{natbib}
%% The amssymb package provides various useful mathematical symbols
\usepackage{amssymb}
%% The amsmath package provides various useful equation environments.
\usepackage{amsmath}
\usepackage{epstopdf}
\usepackage{mathptmx,mathtools}
\usepackage{subcaption}
\usepackage{float} % Add this in your preamble
\usepackage[export]{adjustbox}
\usepackage{array} % Needed for m{} column type
\usepackage[percent]{overpic} % <-- put this in your preamble



\usepackage{amsmath}
\usepackage{nicefrac}
\usepackage{dirtytalk}
\usepackage{graphicx}   % For including images
\usepackage{tikz}
\usetikzlibrary{calc,tikzmark} % needed <<<<<<<<<<
\usepackage{xcolor}     % For coloring
\usepackage{lineno}
\renewcommand\linenumberfont{\normalfont\bfseries\small} % <-- Increase line number size
\usetikzlibrary{arrows.meta}
\usetikzlibrary{calc, positioning}
%%%Author macros
\def\tsc#1{\csdef{#1}{\textsc{\lowercase{#1}}\xspace}}
\tsc{WGM}
\tsc{QE}

\begin{document}
	
	\setcounter{figure}{1}
	
		%%%%%%%%%%%%%%%%%%%%%%%%%%%%%________Figure_3___________%%%%%%%%%%%%%%%%%%%%%
	\begin{figure}[!t]
		\centering
		
		%=================== Subfigure (a) ===================%
		\begin{subfigure}[t]{0.225\textwidth}
			\centering
			\begin{overpic}[width=\textwidth,trim={26.5cm 0cm 24cm 0cm},clip]{Fig3a.eps}
				\put(12,95){\fontfamily{ptm}\textbf{(a) $\Delta\tau_{1-3}$}}
				\put(25,38){\fontfamily{ptm}\textbf{$\tau_1$}}
				\put(25,65){\fontfamily{ptm}\textbf{$\tau_2$}}
				\put(25,89){\fontfamily{ptm}\textbf{$\tau_3$}}
				\put(10,50){\fontfamily{ptm}\textbf{I}}
				\put(10,60){\fontfamily{ptm}\textbf{II}}
				\put(10,76){\fontfamily{ptm}\textbf{III}}
				\put(10,86){\fontfamily{ptm}\textbf{IV}}
			\end{overpic}
		\end{subfigure}
		\hfill
		%=================== Subfigure (b–d) ===================%
		\begin{subfigure}[t]{0.768\textwidth}
			\centering
			\begin{overpic}[width=\textwidth,trim={7cm 0cm 5.5cm 0cm},clip]{Fig3b.eps}
				% Panel labels
				\put(10,61.5){\fontfamily{ptm}\textbf{(b) $\Delta\tau_{1-2}$}}
				\put(44,61.5){\fontfamily{ptm}\textbf{(c) $\Delta\tau_{2-3}$}}
				\put(79,61.5){\fontfamily{ptm}\textbf{(d) $\Delta\tau_{3-4}$}}
				% Region labels (adjust as needed)
				\put(5,15){\fontfamily{ptm}\textbf{I}}
				\put(5,29){\fontfamily{ptm}\textbf{II}}
				\put(5,43){\fontfamily{ptm}\textbf{III}}
				\put(5,57){\fontfamily{ptm}\textbf{IV}}
				
				\put(40,15){\fontfamily{ptm}\textbf{I}}
				\put(40,29){\fontfamily{ptm}\textbf{II}}
				\put(40,43){\fontfamily{ptm}\textbf{III}}
				\put(40,57){\fontfamily{ptm}\textbf{IV}}
				
				\put(74.5,15){\fontfamily{ptm}\textbf{I}}
				\put(74.5,29){\fontfamily{ptm}\textbf{II}}
				\put(74.5,43){\fontfamily{ptm}\textbf{III}}
				\put(74.5,57){\fontfamily{ptm}\textbf{IV}}
			\end{overpic}
		\end{subfigure}
		
		%=================== Caption ===================%
		\caption{Backward FTLE fields of the rising bubble case $C_6$ (a) over the full interval $\Delta\tau_{1-3} = 4 - 20$, with observation windows indicated as I–IV. Temporal evolution of the FTLE fields within these windows, resolved into three equal sub-intervals: (b) $\Delta\tau_{1-2} = 4 - 8$ and (c) $\Delta\tau_{2-3} = 12 - 20$ , and (d) $\Delta\tau_{3-4} = 20 - 28$, showing the temporal evolution of the FTLE field.}
		\label{Fig3}
	\end{figure}
	%%%%%%%%%%%%%%%%%%%%%%%%%%%%%________Figure_3_end___________%%%%%%%%%%%%%%%%%%%%%
	
\end{document}